\documentclass[11pt]{article}

\usepackage{amsmath,amssymb}
\usepackage{xurl}
\usepackage[hidelinks]{hyperref}
\setlength{\emergencystretch}{2em}

% Shared commands for notes in docs/paper-gaps/.
% Notation follows the LDT paper (references/ldt-paper/).

\newcommand{\Lean}{\textsc{Lean}}
\newcommand{\leanid}[1]{\textsf{#1}}
\newcommand{\ghissue}[1]{\href{https://github.com/LionSR/MIPStarRE/issues/#1}{GitHub issue \##1}}
\newcommand{\ghpr}[1]{\href{https://github.com/LionSR/MIPStarRE/pull/#1}{GitHub PR \##1}}

% --- Paper-compatible notation ---
\newcommand{\F}{\mathbb{F}}
\newcommand{\N}{\mathbb{N}}
\newcommand{\C}{\mathbb{C}}
\newcommand{\tr}{\operatorname{tr}}
\newcommand{\ot}{\otimes}
\newcommand{\E}{\mathop{\bf E\/}}
\newcommand{\bx}{\mathbf{x}}
\newcommand{\by}{\mathbf{y}}
\newcommand{\bu}{\mathbf{u}}
\newcommand{\bv}{\mathbf{v}}

% Approximation relation (paper uses \approx_\eps and \simeq_\zeta)
\newcommand{\approxeps}{\approx_{\varepsilon}}
\newcommand{\simgeq}{\simeq_{\zeta}}

% Polynomial spaces (paper uses \polyfunc, \polysub, \polymeas)
\newcommand{\polyfunc}[3]{\operatorname{Poly}(#1,#2,#3)}
\newcommand{\polysub}[3]{\operatorname{PolySub}(#1,#2,#3)}
\newcommand{\polymeas}[3]{\operatorname{PolyMeas}(#1,#2,#3)}


\title{Chapter~3 Same-Space Specializations of Two-Space Paper Propositions}
\author{}
\date{2026-06-11}

\begin{document}
\maketitle

\section{At a glance}

\begin{itemize}
  \item \textbf{Difficulty: medium as a cross-reference audit issue, low as a
  mathematical obstruction.}  No new proof machinery is needed.  The faithful
  two-space forms are already proved, \verb|sorry|-free, and consumed on the
  main-theorem path.
  \item \textbf{Estimated weight:} blueprint-only label changes for five
  propositions; one note-only entry (\verb|consSubMeas|) where the same-space
  form is downstream-sufficient.
  \item \textbf{Mathlib/project split:} 100\% project-local.  The issue
  concerns the chapter~3 blueprint \verb|\lean{}| cross-references.
  \item \textbf{Key mathematical inputs:} the \verb|*_heterogeneous| siblings
  in
  \path{MIPStarRE/LDT/Preliminaries/Triangles/Core.lean},
  \path{.../Triangles/SimEq.lean},
  \path{.../ComparisonCore.lean}, and
  \path{.../ComparisonProjective.lean}.  All relevant declarations are
  project-local; no external Mathlib inputs beyond the ambient linear-algebra
  and finite-dimensionality libraries.
\end{itemize}

\section{Key theorem forms}

The formalization isolates the following mathematical statements, all
concerning the paper's approximation relation \(\simeq_\delta\) between
submeasurements on two Hilbert spaces \cite{Ji2020LowIndividualDegree}.%
\footnote{Tracked in \ghissue{2338}.}

\begin{enumerate}
  \item \textbf{Source theorem form.}  For \(A=\{A_a^x\}\) on \(\mathcal H_A\)
  and \(B=\{B_a^x\}\) on \(\mathcal H_B\) and a state
  \(\ket{\psi}\in\mathcal H_A\otimes\mathcal H_B\), the paper's
  \(\simeq_\delta\) relation is defined by
  \[
    A_a^x\otimes I \simeq_\delta I\otimes B_a^x
    \quad\text{when}\quad
    \E_{x\sim\mathcal D}\sum_{a\neq b}
    \bra{\psi}A_a^x\otimes B_b^x\ket{\psi}\le\delta .
    \tag{\texttt{def:simeq}}
  \]
  The five propositions named below each assert a bound or implication
  involving this two-space relation (see
  \path{references/ldt-paper/preliminaries.tex}).

  \item \textbf{Faithful formal form.}  Each proposition is formalized as a
  \verb|*_heterogeneous| declaration carrying independent index types
  \(\iota_A,\iota_B\) for the left and right submeasurement families and a
  state \(\ket{\psi}\in\mathcal H_{\iota_A}\otimes\mathcal H_{\iota_B}\).%
  \footnote{The declarations are \leanid{triangleSub\_heterogeneous},
  \leanid{simeqTriangleInequality\_heterogeneous},
  \leanid{simeqToApprox\_heterogeneous},
  \leanid{approxToSimeq\_heterogeneous}, and
  \leanid{simeqDataProcessing\_heterogeneous} in
  \doclink{MIPStarRE/LDT/Preliminaries/}.}
  The same-space \((\iota\times\iota)\) specializations also exist and are
  retained for internal use, but the blueprint \verb|\lean{}| links now target
  the two-space forms.

  \item \textbf{Downstream theorem form.}  The main-theorem path consumes the
  \verb|*_heterogeneous| forms at
  \path{MIPStarRE/LDT/Test/MainTheorem/SourceRoleRegister/Final.lean}
  lines~449,~495,~505,~512,~563;
  \path{.../Completion.lean} line~925; and
  \path{.../Core.lean} line~312.  No new hypotheses are needed beyond those
  already in the two-space forms.
\end{enumerate}

The same-space forms are retained as auxiliary internal interfaces (they are
mathematically correct as specializations and are occasionally convenient for
proofs that already operate on a symmetrized state).  However, the
\emph{faithful formalization} of each paper proposition is the
\verb|_heterogeneous| form, which is the version consumed on the main-theorem
path.

\section{Five retargeted links}

The retargeting replaces five blueprint \verb|\lean{}| cross-references.  In each
case the same-space specialization sets \(\iota_A = \iota_B = \iota\) in the
heterogeneous statement, reducing the two-space form
\[
  \ket{\psi}\in\mathcal H_{\iota_A}\otimes\mathcal H_{\iota_B},\qquad
  A_{\iota_A}\otimes I \simeq_\delta I\otimes B_{\iota_B}
\]
to its diagonal
\[
  \ket{\psi}\in\mathcal H_\iota\otimes\mathcal H_\iota,\qquad
  A_\iota\otimes I \simeq_\delta I\otimes B_\iota .
\]
The heterogeneous form is the faithful transcription of
\path{def:simeq} in \cite[Section~3]{Ji2020LowIndividualDegree};
the same-space form is a correct specialization retained as an auxiliary
interface.%

\footnote{The five retargetings map
\leanid{triangleSub}~$\to$~\leanid{triangleSub\_heterogeneous}
(blueprint \verb|prop:triangle-sub|, bound \(\delta+\sqrt\eps\)),
\leanid{simeqTriangleInequality}~$\to$~\leanid{simeqTriangleInequality\_heterogeneous}
(\verb|prop:simeq-triangle-inequality|, bound \(\eps+2\sqrt{\delta+\gamma}\)),
\leanid{simeqToApprox}~$\to$~\leanid{simeqToApprox\_heterogeneous}
(\verb|prop:simeq-to-approx| forward, \(\simeq_\delta\Rightarrow\approx_{2\delta}\)),
\leanid{approxToSimeq}~$\to$~\leanid{approxToSimeq\_heterogeneous}
(converse, projective), and
\leanid{simeqDataProcessing}~$\to$~\leanid{simeqDataProcessing\_heterogeneous}
(\verb|prop:simeq-data-processing|, bound \(\delta\)).
All five \verb|_heterogeneous| forms are \verb|sorry|-free.
Four are consumed on the main-theorem path at
\path{MIPStarRE/LDT/Test/MainTheorem/SourceRoleRegister/Final.lean}
lines~449,~495,~505,~512,~563;
\path{.../Completion.lean} line~925; and
\path{.../Core.lean} line~312.
\leanid{simeqDataProcessing\_heterogeneous} is proved but not yet
consumed on the main-theorem path.}

The bounds match the paper exactly: \(\delta+\sqrt\eps\) (triangle substitution),
\(\eps+2\sqrt{\delta+\gamma}\) (triangle inequality),
\(2\delta\) and \(\delta\) (forward/converse), and \(\delta\) (data
processing).

\section{\texorpdfstring{\leanid{consSubMeas}}{consSubMeas}: a single-case note}

The lemma \verb|prop:cons-sub-meas| (\path{ch03\_preliminaries.tex}:476) is
currently linked to \leanid{consSubMeas}, which is a same-space
(\(\iota\times\iota\)) specialization.  No \verb|consSubMeas\_heterogeneous|
exists, and the submeasurement-family infrastructure that it uses
(\leanid{diagonalSandwichFamily}, \leanid{totalSandwichFamily}) is
\(\iota\times\iota\)-only.

However, the only consumer of \verb|consSubMeas| on the main-theorem path is
\path{MainChain.lean}:209, which applies it at the symmetrized state
\leanid{strategy.state : QuantumState (ι×ι)}.  The global proof never needs
the fully general two-space form of this lemma.  The same-space
\verb|consSubMeas| is therefore \emph{downstream-sufficient} for the current
formal proof of the main theorem.

If a future development path requires \verb|consSubMeas| over \(\iota_A \times
\iota_B\), the infrastructure for constructing the needed diagonal and
sandwich families
(\leanid{mkLeftPlacedSubMeas}/\leanid{mkRightPlacedSubMeas} in
\path{SubMeasurementFamilies.lean}:500--560) already exists and can be used to
prove a \verb|consSubMeas\_heterogeneous| sibling.

\section{Conclusion}

\begin{itemize}
  \item \textbf{Resolved (blueprint links retargeted):} the five
  \verb|_heterogeneous| link retargetings listed above.
  \item \textbf{Documented (same-space form retained as auxiliary):}
  \leanid{triangleSub}, \leanid{simeqTriangleInequality},
  \leanid{simeqToApprox}, \leanid{approxToSimeq},
  \leanid{simeqDataProcessing}, and \leanid{consSubMeas} are retained as
  auxiliary internal interfaces.  They are not removed or renamed.  The
  \verb|_heterogeneous| forms carry the main-theorem load.
  \item \textbf{Deferred:} a \leanid{consSubMeas\_heterogeneous} sibling is not
  currently needed.  If it becomes needed, the submeasurement-family
  infrastructure in \path{SubMeasurementFamilies.lean} supports its
  construction.
\end{itemize}

There is no mathematical discrepancy.  The retargeting makes the blueprint
cross-references more faithful to the paper's two-space definition of
\(\simeq_\delta\).%
\footnote{Related: \ghissue{2337} (same-space cleanup of the \verb|mainFormal|
interface, the Preliminaries-layer analogue).  The same-space forms here are
distinct from the genuine open mathematical gap in \ghissue{1093} (the
\S9~\(\eta\) / total-overlap term).}

\bibliographystyle{alpha}
\bibliography{references}

\end{document}
